\documentclass[11pt,a4paper]{report}

\usepackage[T1]{fontenc}
\usepackage[utf8]{inputenc}
\usepackage[margin=2.6cm]{geometry}
\usepackage{lmodern}
\usepackage{microtype}
\usepackage{longtable}
\usepackage{booktabs}
\usepackage{array}
\usepackage{xcolor}
\usepackage{listings}
\usepackage{makeidx}
\usepackage{enumitem}
\usepackage{fancyhdr}
\usepackage[hidelinks]{hyperref}

\makeindex

% A straight double quote, so that script fragments can be quoted in running
% text without the usual LaTeX curly substitution.
\newcommand{\dq}{\char`\"}

\definecolor{codebg}{gray}{0.96}
\definecolor{codekey}{rgb}{0.10,0.30,0.60}
\definecolor{codestr}{rgb}{0.55,0.20,0.10}
\definecolor{codecom}{rgb}{0.35,0.45,0.35}

\lstdefinestyle{shell}{
  basicstyle=\ttfamily\small,
  backgroundcolor=\color{codebg},
  frame=single,
  framerule=0pt,
  xleftmargin=4pt,
  framexleftmargin=4pt,
  breaklines=true,
  columns=fullflexible,
  keepspaces=true,
  showstringspaces=false,
}

\lstdefinelanguage{scriptjson}{
  basicstyle=\ttfamily\small,
  backgroundcolor=\color{codebg},
  frame=single,
  framerule=0pt,
  xleftmargin=4pt,
  framexleftmargin=4pt,
  breaklines=true,
  columns=fullflexible,
  keepspaces=true,
  showstringspaces=false,
  morestring=[b]",
  stringstyle=\color{codestr},
  comment=[l]{//},
  commentstyle=\color{codecom}\itshape,
}

\lstset{style=shell}

% Argument tables in the command reference.
\newenvironment{argtable}
  {\par\vspace{2pt}\footnotesize
   \setlength{\tabcolsep}{4pt}%
   \begin{longtable}{@{}>{\raggedright\arraybackslash}p{0.28\textwidth}%
                       >{\raggedright\arraybackslash}p{0.15\textwidth}%
                       >{\raggedright\arraybackslash}p{0.20\textwidth}%
                       >{\raggedright\arraybackslash}p{0.27\textwidth}@{}}
   \toprule
   \textbf{Argument} & \textbf{Value} & \textbf{Default} & \textbf{Notes} \\
   \midrule
   \endfirsthead
   \toprule
   \textbf{Argument} & \textbf{Value} & \textbf{Default} & \textbf{Notes} \\
   \midrule
   \endhead
   \bottomrule
   \endfoot}
  {\end{longtable}\normalsize\par}

\pagestyle{fancy}
\fancyhf{}
\fancyhead[L]{\small\leftmark}
\fancyhead[R]{\small\thepage}
\renewcommand{\headrulewidth}{0.4pt}

\title{\Huge SedSAT3 \\[0.3cm]
       \Large Command Line User's Manual \\[0.6cm]
       \large \texttt{sedsat3-cli}}
\author{Sediment Source Assessment Tool, version 1.1.6}
\date{\today}

\begin{document}

\maketitle

\begin{abstract}
\noindent
\texttt{sedsat3-cli} runs SedSAT3 analyses from a script instead of from the
graphical interface. It loads a project, carries out a list of commands against
it in order, and writes a report of what each one produced.

It is the same software. The command line build drives the same analysis code
the windowed application drives, using the same command names and the same
argument names, so anything that can be done through the dialogs can be
written down in a script and repeated exactly.

This manual is for someone running analyses. It assumes familiarity with
sediment fingerprinting and with SedSAT3 itself, but not with programming.
\end{abstract}

\tableofcontents

%======================================================================
\chapter{Introduction}

\section{What it is for}

The graphical interface is the right way to explore a dataset: you can see the
plots, adjust a setting, and look again. It is a poor way to do the same thing
twenty times, and it leaves no record of what was done.

The command line build addresses three things the interface cannot:

\begin{description}[style=nextline,leftmargin=0pt]
  \item[A run you can repeat] A script is a file. It states which project was
    loaded, which constituents were included, and which analyses were carried
    out with which settings. Running it again a year later produces the same
    thing, and anyone else can run it too.
  \item[A run you can vary] Comparing five tracer sets means five scripts
    differing in one line, not five sessions of clicking and remembering.
  \item[A run you can automate] It needs no display, so it can run on a server,
    in a scheduled job, or as part of a test that checks the results have not
    changed.
\end{description}

\section{What it is not}

It does not replace the interface. There are no plots, and there is no way to
look at an intermediate result and change your mind. Use the interface to
decide what to do, and a script to do it repeatedly.

One limitation is worth stating plainly at the outset: \textbf{the command line
build cannot import from Excel}. It requires a project file that has already
been saved from the interface. Importing a spreadsheet involves choosing which
sheets to use and what each one holds, and those choices are made in dialogs
that have no command line equivalent yet.

\section{How this manual is arranged}

Chapter~\ref{ch:installing} covers obtaining and building the program.
Chapter~\ref{ch:quickstart} walks through a first run from beginning to end.
Chapter~\ref{ch:scripts} describes the script file in detail.
Chapter~\ref{ch:setup} covers the commands that configure a project, and
chapter~\ref{ch:commands} is the reference for every analysis command.
Chapter~\ref{ch:report} explains the report that a run produces, and
chapter~\ref{ch:troubleshooting} covers what to do when something goes wrong.

%======================================================================
\chapter{Installing}
\label{ch:installing}

\section{Requirements}

The command line build needs the same libraries as the graphical one: Qt~6,
the GNU Scientific Library, Armadillo, and a C++17 compiler. It links Qt
Widgets but never opens a window, so it runs on a machine with no display.

\section{Building}

From the top of the source tree:

\begin{lstlisting}
qmake6 SedSat3CLI.pro
make -j8
\end{lstlisting}

This produces \texttt{sedsat3-cli} in the build directory. On Windows, use the
\texttt{qmake} supplied with the Qt~6 kit and the build tool that goes with
your compiler; the project file is the same.

\begin{quote}
\small
On systems with both Qt~5 and Qt~6 installed, the plain \texttt{qmake} command
often resolves to Qt~5 and the build fails immediately. Use \texttt{qmake6},
or the full path to the Qt~6 one.
\end{quote}

\section{Telling it where the resources are}

The program reads \texttt{resources/forms\_structures.json}, which defines
every command and its arguments. It uses that file to check a script before
running it and to answer \texttt{--list-commands} and \texttt{--describe}.

It looks for the file next to the program and in the usual places relative to
it. If your build directory is somewhere unusual, say where to look:

\begin{lstlisting}
export SEDSAT3_RESOURCES=/path/to/SedSat3/resources
\end{lstlisting}

or on Windows:

\begin{lstlisting}
set SEDSAT3_RESOURCES=C:\path\to\SedSat3\resources
\end{lstlisting}

The \texttt{--resources} option does the same thing for a single run.

If the file cannot be found the program still runs, but it warns you, and two
things stop working: scripts are no longer checked, and arguments you leave out
are no longer filled in with their defaults. The second of those matters more
than it sounds, so it is worth setting the variable properly.

%======================================================================
\chapter{A first run}
\label{ch:quickstart}

\section{Seeing what is available}

\begin{lstlisting}
sedsat3-cli --list-commands
\end{lstlisting}

lists every command, and

\begin{lstlisting}
sedsat3-cli --describe "Levenberg-Marquardt"
\end{lstlisting}

prints one command's description and its arguments, with the default for each.
These come from the same definitions the dialogs are built from, so they
describe the software as it actually is.

\section{Writing a script}

A script is a text file containing JSON. Save the following as
\texttt{run.json}, next to your project file:

\begin{lstlisting}[language=scriptjson]
{
  "project": "my-project.cmb",
  "output": "report.json",
  "steps": [
    {
      "command": "Levenberg-Marquardt",
      "arguments": {
        "Sample": "CTAIL1",
        "Apply size and organic matter correction": "false"
      }
    }
  ]
}
\end{lstlisting}

That is the whole thing: the project to load, where to put the report, and one
analysis to carry out.

\section{Running it}

\begin{lstlisting}
sedsat3-cli run.json
\end{lstlisting}

You will see the step announced and its progress, and then where the report was
written:

\begin{lstlisting}
[1] Levenberg-Marquardt
  100%
report written to /home/you/work/report.json
\end{lstlisting}

\section{Reading the result}

The report is JSON. Each step records what was asked for, whether it succeeded,
how long it took, and what it produced. The source contributions from the run
above are under that step's results:

\begin{lstlisting}[language=scriptjson]
"1:Contributions": {
  "Bank":    0.2122,
  "Crop":    0.4104,
  "Forest":  0.1822,
  "Pasture": 0.1953
}
\end{lstlisting}

If you have the \texttt{jq} utility, it is a convenient way to pick things out:

\begin{lstlisting}
jq '.steps[] | {command, status, seconds}' report.json
\end{lstlisting}

%======================================================================
\chapter{Scripts}
\label{ch:scripts}

\section{Structure}

A script is a JSON object with four keys. Only two of them are required.

\begin{longtable}{@{}p{0.22\textwidth}p{0.70\textwidth}@{}}
\toprule
\textbf{Key} & \textbf{Meaning} \\
\midrule
\endfirsthead
\texttt{project} & The project file to load. Required. \\
\addlinespace
\texttt{steps} & The commands to carry out, as an array. Required. \\
\addlinespace
\texttt{output} & Where to write the report. If absent, the report is printed
  to standard output. \\
\addlinespace
\texttt{on\_error} & \texttt{\dq{}stop\dq{}} (the default) abandons the run
  after a step fails; \texttt{\dq{}continue\dq{}} carries on with the
  remaining steps. \\
\addlinespace
\texttt{working\_folder} & Where analyses that write files of their own should
  put them. Defaults to the directory holding the script. \\
\bottomrule
\end{longtable}

\section{Steps run in order, and they share one project}

\texttt{steps} is an array rather than an object, and that is deliberate: the
order matters. All the steps act on one dataset, and several of them change
it. Outlier removal, the organic matter and particle size correction, and the
element selection commands all alter what the steps after them will see.

This is the main thing to understand about scripts. A script is not a list of
independent jobs; it is a sequence, in the same way that a session in the
interface is a sequence.

\section{Steps}

Each step is an object with a command and its arguments:

\begin{lstlisting}[language=scriptjson]
{
  "command": "Outlier",
  "arguments": { "Threshold": "3" }
}
\end{lstlisting}

The command names and the argument names are exactly those the dialogs use,
including their capitalisation and spacing. \texttt{--describe} will tell you
what a command accepts.

Argument values are written as text. Numbers and \texttt{true}/\texttt{false}
are also accepted and converted, so both of these work:

\begin{lstlisting}[language=scriptjson]
"arguments": { "Number of samples": "1000", "Dissolve Chains": "true"  }
"arguments": { "Number of samples": 1000,   "Dissolve Chains": true    }
\end{lstlisting}

\section{You need only give what you want to change}

Anything you leave out is filled in with the default the dialog would have
shown. The report lists what was filled in, under \texttt{defaulted}, so the
record of the run is complete even where the script was not.

This is worth relying on. Specify the two or three settings that matter for
what you are doing and leave the rest alone; the run is still reproducible,
because the report records the values that were actually used.

\section{Paths}

A path in a script that is not absolute is interpreted relative to the
directory holding the script. A script and its project can therefore be moved
or copied together, and sent to someone else, without editing.

Files that an analysis writes itself, such as the MCMC sample output, go to
\texttt{working\_folder}, which defaults to the script's directory.

\section{Checking before running}

Every step is checked before any of them runs: that the command exists, and
that each argument is one the command accepts. If anything is wrong, nothing
is carried out and the problems are listed:

\begin{lstlisting}
error: The script was not run because of the following:
  step 1: Unknown command 'Levenberg-Marquart'. Did you mean
          'Levenberg-Marquardt'?
  step 2: Command 'Levenberg-Marquardt' has no argument 'Sampel'.
          Did you mean 'Sample'?
\end{lstlisting}

Checking everything first matters because of the previous section. A script
that stopped halfway through because of a spelling mistake in the last step
would leave the project in a state you did not ask for.

%======================================================================
\chapter{Configuring the project}
\label{ch:setup}

The analyses work on whatever the project file was saved with: which
constituents are included, which samples are in use, what each constituent is,
and which group holds the targets. In the interface those are set in tables and
dialogs rather than through the analysis menus.

The four commands in this chapter set them from a script. Where one of them
disagrees with the project file, the command wins: running one is how a script
says it wants something other than what was saved.

They are what make a script useful for more than replaying a saved
configuration. Comparing tracer sets, or correcting an isotope's settings, is a
matter of one step.

\section{\texttt{set-target-group}}

Chooses which group holds the target samples.

\begin{argtable}
\texttt{group} & name & --- & a group name in the project \\
\end{argtable}

\begin{lstlisting}[language=scriptjson]
{ "command": "set-target-group", "arguments": { "group": "Mixture" } }
\end{lstlisting}

\section{\texttt{set-elements}}

Chooses which constituents take part.

\begin{argtable}
\texttt{only} & name or list & --- & include exactly these, exclude everything else \\
\texttt{include} & name or list & --- & include these, leave the rest alone \\
\texttt{exclude} & name or list & --- & exclude these, leave the rest alone \\
\end{argtable}

\texttt{only} is applied first, then \texttt{include}, then \texttt{exclude}.
Every name is checked before anything is changed, so a misspelling cannot leave
the project half configured.

\begin{lstlisting}[language=scriptjson]
{ "command": "set-elements",
  "arguments": { "only": ["Aluminum", "Barium", "Lead", "Zinc"] } }
\end{lstlisting}

\section{\texttt{set-samples}}

Chooses which samples in a group take part.

\begin{argtable}
\texttt{group} & name & --- & required \\
\texttt{only} & name or list & --- & include exactly these within the group \\
\texttt{include} & name or list & --- & include these \\
\texttt{exclude} & name or list & --- & exclude these \\
\end{argtable}

\begin{lstlisting}[language=scriptjson]
{ "command": "set-samples",
  "arguments": { "group": "Bank", "exclude": ["B7", "B12"] } }
\end{lstlisting}

\section{\texttt{set-element-role}}

Sets what a constituent is, and for an isotope what it is measured against.

\begin{argtable}
\texttt{element} & name & --- & required \\
\texttt{role} & name & --- & \texttt{Element}, \texttt{Isotope},
  \texttt{ParticleSize}, \texttt{OM} or \texttt{DoNotInclude} \\
\texttt{base element} & name & --- & for an isotope: the constituent supplying
  its concentration \\
\texttt{standard ratio} & number & --- & for an isotope: the standard isotopic
  ratio \\
\end{argtable}

\begin{lstlisting}[language=scriptjson]
{ "command": "set-element-role",
  "arguments": { "element": "C13", "role": "Isotope",
                 "base element": "TOC", "standard ratio": "0.011113" } }
\end{lstlisting}

\begin{quote}
\small
\textbf{Isotopes need their base element to be part of the analysis.} An
isotope is measured as a delta value, and converting that to a concentration
requires the concentration of the constituent it is referenced against. If the
base element is not included, the conversion has nothing to work with.

A base element must also have the role \texttt{Element}. Constituents whose
role is \texttt{OM} or \texttt{ParticleSize} are set aside before the analyses
run, so an isotope referenced against organic carbon needs that constituent
given the role \texttt{Element} and included:

\begin{lstlisting}[language=scriptjson]
{ "command": "set-element-role",
  "arguments": { "element": "TOC", "role": "Element" } },
{ "command": "set-elements", "arguments": { "include": ["TOC"] } }
\end{lstlisting}

Note that this also makes the base element a fingerprinting tracer in its own
right, which changes the apportionment. That is a modelling decision, so make
it deliberately.
\end{quote}

%======================================================================
\chapter{Analysis commands}
\label{ch:commands}

Every command the software accepts is listed here, with its arguments, the
values each one takes, and its default. The descriptions are the same ones the
dialogs show.

The \textbf{Value} column says what kind of thing to write in the script.
Where it says ``name'', the accepted names come from the loaded project --- a
target sample, a source group or a constituent --- and the \textbf{Notes}
column says which.

% Generated by docs/manual/generate-command-reference.py
% Do not edit by hand; edit resources/forms_structures.json
% and regenerate, so that the manual and the software agree.

\subsection{\texttt{ANOVA}}
\label{cmd:anova}
\index{ANOVA}

This tool performs analysis of variance. The resulting p-value indicates the power of a certain element in discriminating between all source groups.

\begin{argtable}
\texttt{Box-cox transformation} & \texttt{true} / \texttt{false} & \texttt{\dq{}false\dq{}} & --- \\
\texttt{Log Transformation} & \texttt{true} / \texttt{false} & \texttt{\dq{}false\dq{}} & --- \\
\texttt{Modify the included elements based on the results} & \texttt{true} / \texttt{false} & \texttt{\dq{}false\dq{}} & --- \\
\texttt{OM and Size Correct based on target sample} & name & --- & a target sample name, or empty for all \\
\texttt{P-value threshold} & number or text & \texttt{\dq{}0.05\dq{}} & --- \\
\texttt{Use only selected elements} & \texttt{true} / \texttt{false} & \texttt{\dq{}false\dq{}} & --- \\
\texttt{Use only selected samples} & \texttt{true} / \texttt{false} & \texttt{\dq{}false\dq{}} & --- \\
\end{argtable}

\subsection{\texttt{AutoSelect}}
\label{cmd:autoselect}
\index{AutoSelect}

The \textbf{Auto-select} tool selects a minimum number of elements with discriminatory power to differenciate between all source pair combinations. It performs a \textit{t-test} between each pair of source groups and select the elements with the lowest p-values from each pair. The number of elements choosen from each pair is indicated in the box labeled 'Number of elements from each pair'

\begin{argtable}
\texttt{Include Isotopes} & \texttt{true} / \texttt{false} & \texttt{\dq{}false\dq{}} & --- \\
\texttt{Modify the included elements based on the results} & \texttt{true} / \texttt{false} & \texttt{\dq{}false\dq{}} & --- \\
\texttt{Number of elements from each pair} & whole number & \texttt{\dq{}3\dq{}} & --- \\
\texttt{OM and Size Correct based on target sample} & name & --- & a target sample name, or empty for all \\
\end{argtable}

\subsection{\texttt{BoxCox}}
\label{cmd:boxcox}
\index{BoxCox}

This tool calculates the optimal lambda value in the Box-Cox transformation for the elements in the specified source group to achieve normality.

\begin{argtable}
\texttt{Source/Target group} & name & --- & a source group name \\
\texttt{Use only selected elements} & \texttt{true} / \texttt{false} & \texttt{\dq{}false\dq{}} & --- \\
\texttt{Use only selected samples} & \texttt{true} / \texttt{false} & \texttt{\dq{}false\dq{}} & --- \\
\end{argtable}

\subsection{\texttt{Bracketing Analysis}}
\label{cmd:bracketing-analysis}
\index{Bracketing Analysis}

This tool shows whether the elemental content of the elements in the specified target sample is outside of the range of the elemental contents in all source groups.

\begin{argtable}
\texttt{Correct based on size and organic matter} & \texttt{true} / \texttt{false} & \texttt{\dq{}false\dq{}} & --- \\
\texttt{Sample} & name & --- & a target sample name \\
\texttt{Use only selected elements} & \texttt{true} / \texttt{false} & \texttt{\dq{}false\dq{}} & --- \\
\texttt{Use only selected samples} & \texttt{true} / \texttt{false} & \texttt{\dq{}false\dq{}} & --- \\
\end{argtable}

\subsection{\texttt{Bracketing Analysis Batch}}
\label{cmd:bracketing-analysis-batch}
\index{Bracketing Analysis Batch}

This tool shows whether the elemental content of the elements in target samples is outside of the range of the elemental contents in all source groups for all target samples.

\begin{argtable}
\texttt{Correct based on size and organic matter} & \texttt{true} / \texttt{false} & \texttt{\dq{}false\dq{}} & --- \\
\texttt{Use only selected elements} & \texttt{true} / \texttt{false} & \texttt{\dq{}false\dq{}} & --- \\
\texttt{Use only selected samples} & \texttt{true} / \texttt{false} & \texttt{\dq{}false\dq{}} & --- \\
\end{argtable}

\subsection{\texttt{CMB Bayesian}}
\label{cmd:cmb-bayesian}
\index{CMB Bayesian}

This tool performs Bayesian fingerprinting using Markov-Chain Monte Carlo algorithm for the specified target samle.

\begin{argtable}
\texttt{Apply size and organic matter correction} & \texttt{true} / \texttt{false} & \texttt{\dq{}false\dq{}} & --- \\
\texttt{Dissolve Chains} & \texttt{true} / \texttt{false} & \texttt{\dq{}true\dq{}} & --- \\
\texttt{Number of chains} & whole number & \texttt{\dq{}8\dq{}} & --- \\
\texttt{Number of samples} & whole number & \texttt{\dq{}1000\dq{}} & --- \\
\texttt{Sample} & name & --- & a target sample name \\
\texttt{Samples File Name} & file name & \texttt{\dq{}MCMC\_Output.txt\dq{}} & --- \\
\texttt{Samples to be discarded (burnout)} & whole number & \texttt{\dq{}100\dq{}} & --- \\
\end{argtable}

\subsection{\texttt{CMB Bayesian-Batch}}
\label{cmd:cmb-bayesian-batch}
\index{CMB Bayesian-Batch}

This tool performs Bayesian fingerprinting using Markov-Chain Monte Carlo algorithm for all target samle.

\begin{argtable}
\texttt{Apply size and organic matter correction} & \texttt{true} / \texttt{false} & \texttt{\dq{}false\dq{}} & --- \\
\texttt{Dissolve Chains} & \texttt{true} / \texttt{false} & \texttt{\dq{}true\dq{}} & --- \\
\texttt{Number of chains} & whole number & \texttt{\dq{}8\dq{}} & --- \\
\texttt{Number of samples} & whole number & \texttt{\dq{}1000\dq{}} & --- \\
\texttt{Samples File Name} & file name & \texttt{\dq{}MCMC\_Output.txt\dq{}} & --- \\
\texttt{Samples to be discarded (burnout)} & whole number & \texttt{\dq{}100\dq{}} & --- \\
\end{argtable}

\subsection{\texttt{CorMat}}
\label{cmd:cormat}
\index{CorMat}

This tool generates the element correlation matrix for the chosen source group. To perform the correlation analysis on the corrected elemental profiles based on size and organic matter, select the target sample based on which the correction to be performed.

\begin{argtable}
\texttt{OM and Size Correct based on target sample} & name & --- & a target sample name, or empty for all \\
\texttt{Source/Target group} & name & --- & a source group name \\
\texttt{Threshold} & number or text & \texttt{\dq{}0.75\dq{}} & --- \\
\texttt{Use only selected elements} & \texttt{true} / \texttt{false} & \texttt{\dq{}false\dq{}} & --- \\
\texttt{Use only selected samples} & \texttt{true} / \texttt{false} & \texttt{\dq{}false\dq{}} & --- \\
\end{argtable}

\subsection{\texttt{CovMat}}
\label{cmd:covmat}
\index{CovMat}

This tool calculates the covariance matrix between the elemental contents of various elements for the specified source group.

\begin{argtable}
\texttt{Source/Target group} & name & --- & a source group name \\
\end{argtable}

\subsection{\texttt{DF}}
\label{cmd:df}
\index{DF}

This tool shows fitted normal and log-normal distributions to the elemental content of the specified element in the specified source group.

\begin{argtable}
\texttt{Box-cox transformation} & \texttt{true} / \texttt{false} & \texttt{\dq{}false\dq{}} & --- \\
\texttt{Constituent} & name & --- & a constituent name \\
\texttt{OM and Size Correct based on target sample} & name & --- & a target sample name, or empty for all \\
\texttt{Source/Target group} & name & --- & a source group name \\
\texttt{Use only selected samples} & \texttt{true} / \texttt{false} & \texttt{\dq{}false\dq{}} & --- \\
\end{argtable}

\subsection{\texttt{DFA}}
\label{cmd:dfa}
\index{DFA}

This tool performs Discriminant Function Analysis between the two source groups indicated. It provides the p-value, and the scatter matrix of projected elemental profiles based on the principle eigen vector of the discriminant hyper-plane.

\begin{argtable}
\texttt{Box-cox transformation} & \texttt{true} / \texttt{false} & \texttt{\dq{}false\dq{}} & --- \\
\texttt{OM and Size Correct based on target sample} & name & --- & a target sample name, or empty for all \\
\texttt{Source/Target group I} & name & --- & a source group name \\
\texttt{Source/Target group II} & name & --- & a source group name \\
\texttt{Use only selected elements} & \texttt{true} / \texttt{false} & \texttt{\dq{}false\dq{}} & --- \\
\texttt{Use only selected samples} & \texttt{true} / \texttt{false} & \texttt{\dq{}false\dq{}} & --- \\
\end{argtable}

\subsection{\texttt{DFAM}}
\label{cmd:dfam}
\index{DFAM}

This tool performs step-wise Discriminant Function Analysis between the two source groups indicated. It provides the S-value as a result of including various number of elements into the analysis, starting from the element that produces the maximum seperation.

\begin{argtable}
\texttt{Box-cox transformation} & \texttt{true} / \texttt{false} & \texttt{\dq{}false\dq{}} & --- \\
\texttt{OM and Size Correct based on target sample} & name & --- & a target sample name, or empty for all \\
\texttt{Use only selected elements} & \texttt{true} / \texttt{false} & \texttt{\dq{}false\dq{}} & --- \\
\texttt{Use only selected samples} & \texttt{true} / \texttt{false} & \texttt{\dq{}false\dq{}} & --- \\
\end{argtable}

\subsection{\texttt{DFAOnevsRest}}
\label{cmd:dfaonevsrest}
\index{DFAOnevsRest}

This tool performs step-wise Discriminant Function Analysis between the two source groups indicated. It provides the S-value as a result of including various number of elements into the analysis, starting from the element that produces the maximum seperation.

\begin{argtable}
\texttt{Box-cox transformation} & \texttt{true} / \texttt{false} & \texttt{\dq{}false\dq{}} & --- \\
\texttt{OM and Size Correct based on target sample} & name & --- & a target sample name, or empty for all \\
\texttt{Source group} & name & --- & a source group name \\
\texttt{Use only selected elements} & \texttt{true} / \texttt{false} & \texttt{\dq{}false\dq{}} & --- \\
\texttt{Use only selected samples} & \texttt{true} / \texttt{false} & \texttt{\dq{}false\dq{}} & --- \\
\end{argtable}

\subsection{\texttt{EDP}}
\label{cmd:edp}
\index{EDP}

This tool calculates the measures of discriminant power of elements between the two specified source groups. The measures include within and between group difference ratio, discriminant fraction, and the T-test p-value.

\begin{argtable}
\texttt{Box-cox transformation} & \texttt{true} / \texttt{false} & \texttt{\dq{}false\dq{}} & --- \\
\texttt{OM and Size Correct based on target sample} & name & --- & a target sample name, or empty for all \\
\texttt{P-value threshold} & number or text & \texttt{\dq{}0.05\dq{}} & --- \\
\texttt{Source/Target group I} & name & --- & a source group name \\
\texttt{Source/Target group II} & name & --- & a source group name \\
\texttt{Use only selected samples} & \texttt{true} / \texttt{false} & \texttt{\dq{}false\dq{}} & --- \\
\end{argtable}

\subsection{\texttt{EDPM}}
\label{cmd:edpm}
\index{EDPM}

This tool calculates the measures of discriminant power of elements between all pairs of source groups. The measures include within and between group difference ratio, discriminant fraction, and the T-test p-value.

\begin{argtable}
\texttt{Box-cox transformation} & \texttt{true} / \texttt{false} & \texttt{\dq{}false\dq{}} & --- \\
\texttt{Include target samples} & \texttt{true} / \texttt{false} & \texttt{\dq{}false\dq{}} & --- \\
\texttt{OM and Size Correct based on target sample} & name & --- & a target sample name, or empty for all \\
\texttt{P-value threshold} & number or text & \texttt{\dq{}0.05\dq{}} & --- \\
\texttt{Use only selected samples} & \texttt{true} / \texttt{false} & \texttt{\dq{}false\dq{}} & --- \\
\end{argtable}

\subsection{\texttt{Error\_Analysis}}
\label{cmd:error_analysis}
\index{Error\_Analysis}

This tool uses Monte-Carlo simulation via bootstrapping the evaluate the effect of elmination of the indicated percentage of source samples on the fingerprinting results.

\begin{argtable}
\texttt{Apply size and organic matter correction} & \texttt{true} / \texttt{false} & \texttt{\dq{}true\dq{}} & --- \\
\texttt{Number of realizations} & whole number & \texttt{\dq{}1000\dq{}} & --- \\
\texttt{Pecentage eliminated} & whole number & \texttt{\dq{}20\dq{}} & --- \\
\texttt{Sample} & name & --- & a target sample name \\
\texttt{Softmax transformation} & \texttt{true} / \texttt{false} & \texttt{\dq{}true\dq{}} & --- \\
\end{argtable}

\subsection{\texttt{GA}}
\label{cmd:ga}
\index{GA}

This tool performs sediment fingerprinting using the \textbf{Genetic Algorithm} optimization. The mean and standard deviations describing the true distribution of the elemental profiles of each element for each source are estimated as part of the maximization of the likelihood of observing the measured target sample elemental profile.

\begin{argtable}
\texttt{Apply size and organic matter correction} & \texttt{true} / \texttt{false} & \texttt{\dq{}false\dq{}} & --- \\
\texttt{Cross-over probability} & number or text & \texttt{\dq{}1.0\dq{}} & --- \\
\texttt{GA output file} & file name & \texttt{\dq{}GA\_Output.txt\dq{}} & --- \\
\texttt{Mutation probability} & number or text & \texttt{\dq{}0.02\dq{}} & --- \\
\texttt{Number of Generations} & whole number & \texttt{\dq{}100\dq{}} & --- \\
\texttt{Number of threads to be used} & whole number & \texttt{\dq{}8\dq{}} & --- \\
\texttt{Numerical cross-over} & \texttt{true} / \texttt{false} & \texttt{\dq{}false\dq{}} & --- \\
\texttt{Population} & whole number & \texttt{\dq{}100\dq{}} & --- \\
\texttt{Sample} & name & --- & a target sample name \\
\texttt{Shake coefficient} & number or text & \texttt{\dq{}0.05\dq{}} & --- \\
\texttt{Shake coefficient reduction factor} & number or text & \texttt{\dq{}0.75\dq{}} & --- \\
\texttt{Steepest Descent} & \texttt{true} / \texttt{false} & \texttt{\dq{}true\dq{}} & --- \\
\end{argtable}

\subsection{\texttt{GA (disregarding targets)}}
\label{cmd:ga-disregarding-targets}
\index{GA (disregarding targets)}

\begin{argtable}
\texttt{Cross-over probability} & number or text & \texttt{\dq{}1.0\dq{}} & --- \\
\texttt{GA output file} & file name & \texttt{\dq{}GA\_Output.txt\dq{}} & --- \\
\texttt{Mutation probability} & number or text & \texttt{\dq{}0.02\dq{}} & --- \\
\texttt{Number of Generations} & whole number & \texttt{\dq{}40\dq{}} & --- \\
\texttt{Number of threads to be used} & whole number & \texttt{\dq{}8\dq{}} & --- \\
\texttt{Numerical cross-over} & \texttt{true} / \texttt{false} & \texttt{\dq{}false\dq{}} & --- \\
\texttt{Population} & whole number & \texttt{\dq{}40\dq{}} & --- \\
\texttt{Shake coefficient} & number or text & \texttt{\dq{}0.05\dq{}} & --- \\
\texttt{Shake coefficient reduction factor} & number or text & \texttt{\dq{}0.75\dq{}} & --- \\
\end{argtable}

\subsection{\texttt{GA (fixed elemental contribution)}}
\label{cmd:ga-fixed-elemental-contribution}
\index{GA (fixed elemental contribution)}

This tool performs sediment fingerprinting using the \textbf{Genetic Algorithm} optimization. The mean and standard deviations describing the true distribution of the elemental profiles of each element for each source are estimated by direct calculation based on the samples.

\begin{argtable}
\texttt{Apply size and organic matter correction} & \texttt{true} / \texttt{false} & \texttt{\dq{}false\dq{}} & --- \\
\texttt{Cross-over probability} & number or text & \texttt{\dq{}1.0\dq{}} & --- \\
\texttt{GA output file} & file name & \texttt{\dq{}GA\_Output.txt\dq{}} & --- \\
\texttt{Mutation probability} & number or text & \texttt{\dq{}0.02\dq{}} & --- \\
\texttt{Number of Generations} & whole number & \texttt{\dq{}40\dq{}} & --- \\
\texttt{Number of threads to be used} & whole number & \texttt{\dq{}8\dq{}} & --- \\
\texttt{Numerical cross-over} & \texttt{true} / \texttt{false} & \texttt{\dq{}false\dq{}} & --- \\
\texttt{Population} & whole number & \texttt{\dq{}40\dq{}} & --- \\
\texttt{Sample} & name & --- & a target sample name \\
\texttt{Shake coefficient} & number or text & \texttt{\dq{}0.05\dq{}} & --- \\
\texttt{Shake coefficient reduction factor} & number or text & \texttt{\dq{}0.75\dq{}} & --- \\
\end{argtable}

\subsection{\texttt{KS}}
\label{cmd:ks}
\index{KS}

This tool provides visualization of Kolomogrov-Smirnov normality test for the indicated source group. It can be used to evaluate how close the elemental content of all elements are to normal or log-normal distribution for the specified source group.

\begin{argtable}
\texttt{Distribution} & name & --- & one of \texttt{Normal}, \texttt{Lognormal} \\
\texttt{Source/Target group} & name & --- & a source group name \\
\end{argtable}

\subsection{\texttt{KS-individual}}
\label{cmd:ks-individual}
\index{KS-individual}

This tool provides visualization of Kolomogrov-Smirnov normality test for the indicated source group. It can be used to evaluate how close the elemental content of a given element is to normal or log-normal distribution for the specified source group.

\begin{argtable}
\texttt{Constituent} & name & --- & a constituent name \\
\texttt{Distribution} & name & --- & one of \texttt{Normal}, \texttt{Lognormal} \\
\texttt{Source/Target group} & name & --- & a source group name \\
\end{argtable}

\subsection{\texttt{Levenberg-Marquardt}}
\label{cmd:levenberg-marquardt}
\index{Levenberg-Marquardt}

This tool uses the Levenberg-Marquardt optimization algorithm to perform sediment fingerprinting by minimizing the RMSE between the modeled and observed elemental profiles of the target sample. The mean and standard deviations representing the true distribution of elemental profiles for each element and source are directly calculated from the samples.

\begin{argtable}
\texttt{Apply size and organic matter correction} & \texttt{true} / \texttt{false} & \texttt{\dq{}true\dq{}} & --- \\
\texttt{Sample} & name & --- & a target sample name \\
\texttt{Softmax transformation} & \texttt{true} / \texttt{false} & \texttt{\dq{}true\dq{}} & --- \\
\end{argtable}

\subsection{\texttt{Levenberg-Marquardt-Batch}}
\label{cmd:levenberg-marquardt-batch}
\index{Levenberg-Marquardt-Batch}

This tool performs sediment fingerprinting using the \textbf{Levenberg-Marquardt} optimization algorithm on all target samples. The mean and standard deviations describing the true distribution of the elemental profiles of each element for each source are estimated by direct calculation based on the samples.This tool uses the Levenberg-Marquardt optimization algorithm to perform sediment fingerprinting by minimizing the RMSE between the modeled and observed elemental profiles of the target sample. The mean and standard deviations representing the true distribution of elemental profiles for each element and source are directly calculated from the samples.

\begin{argtable}
\texttt{Apply size and organic matter correction} & \texttt{true} / \texttt{false} & \texttt{\dq{}true\dq{}} & --- \\
\texttt{Softmax transformation} & \texttt{true} / \texttt{false} & \texttt{\dq{}true\dq{}} & --- \\
\end{argtable}

\subsection{\texttt{MLR}}
\label{cmd:mlr}
\index{MLR}

This tool performs the regression analysis for organic and size correction on all source groups.

\begin{argtable}
\texttt{Equation} & name & --- & one of \texttt{Linear}, \texttt{Power} \\
\texttt{Organic Matter constituent} & name & --- & a constituent name, or empty, or the organic matter constituent \\
\texttt{P-value threshold} & number or text & \texttt{\dq{}0.05\dq{}} & --- \\
\texttt{Particle Size constituent} & name & --- & a constituent name, or empty, or the particle size constituent \\
\texttt{Use only selected samples} & \texttt{true} / \texttt{false} & \texttt{\dq{}false\dq{}} & --- \\
\end{argtable}

\subsection{\texttt{OM-Size Correct}}
\label{cmd:om-size-correct}
\index{OM-Size Correct}

This tool shows the corrected elemental profiles of source groups after adjusting based on size and organic matter. The correction is done based on the size and organic matter associated with the selected target. The regressions as a result of the last organic and size correction will be used to perform the adjustment.

\begin{argtable}
\texttt{Sample} & name & --- & a target sample name \\
\texttt{Use only selected elements} & \texttt{true} / \texttt{false} & \texttt{\dq{}false\dq{}} & --- \\
\texttt{Use only selected samples} & \texttt{true} / \texttt{false} & \texttt{\dq{}false\dq{}} & --- \\
\end{argtable}

\subsection{\texttt{Outlier}}
\label{cmd:outlier}
\index{Outlier}

This tool conducts outlier analysis on the elemental profiles for the specified source group. The analysis begins by applying a Box-Cox transformation to each elemental content, followed by calculating the standardized value using the formula (x - mean) / std. Any absolute values exceeding the user-defined threshold are highlighted.

\begin{argtable}
\texttt{Source/Target group} & name & --- & a source group name \\
\texttt{Threshold} & number or text & \texttt{\dq{}3\dq{}} & --- \\
\texttt{Use only selected elements} & \texttt{true} / \texttt{false} & \texttt{\dq{}false\dq{}} & --- \\
\texttt{Use only selected samples} & \texttt{true} / \texttt{false} & \texttt{\dq{}false\dq{}} & --- \\
\end{argtable}

\subsection{\texttt{SDFA}}
\label{cmd:sdfa}
\index{SDFA}

This tool performs step-wise Discriminant Function Analysis between the two source groups indicated. It provides the S-value as a result of including various number of elements into the analysis, starting from the element that produces the maximum seperation.

\begin{argtable}
\texttt{Box-cox transformation} & \texttt{true} / \texttt{false} & \texttt{\dq{}false\dq{}} & --- \\
\texttt{OM and Size Correct based on target sample} & name & --- & a target sample name, or empty for all \\
\texttt{Source/Target group I} & name & --- & a source group name \\
\texttt{Source/Target group II} & name & --- & a source group name \\
\texttt{Use only selected elements} & \texttt{true} / \texttt{false} & \texttt{\dq{}false\dq{}} & --- \\
\texttt{Use only selected samples} & \texttt{true} / \texttt{false} & \texttt{\dq{}false\dq{}} & --- \\
\end{argtable}

\subsection{\texttt{SDFAM}}
\label{cmd:sdfam}
\index{SDFAM}

This tool performs step-wise Discriminant Function Analysis between the two source groups indicated. It provides the S-value as a result of including various number of elements into the analysis, starting from the element that produces the maximum seperation.

\begin{argtable}
\texttt{Box-cox transformation} & \texttt{true} / \texttt{false} & \texttt{\dq{}false\dq{}} & --- \\
\texttt{Modify the included elements based on the results} & \texttt{true} / \texttt{false} & \texttt{\dq{}false\dq{}} & --- \\
\texttt{OM and Size Correct based on target sample} & name & --- & a target sample name, or empty for all \\
\texttt{Use only selected elements} & \texttt{true} / \texttt{false} & \texttt{\dq{}false\dq{}} & --- \\
\texttt{Use only selected samples} & \texttt{true} / \texttt{false} & \texttt{\dq{}false\dq{}} & --- \\
\end{argtable}

\subsection{\texttt{SDFAOnevsRest}}
\label{cmd:sdfaonevsrest}
\index{SDFAOnevsRest}

This tool performs step-wise Discriminant Function Analysis between the two source groups indicated. It provides the S-value as a result of including various number of elements into the analysis, starting from the element that produces the maximum seperation.

\begin{argtable}
\texttt{Box-cox transformation} & \texttt{true} / \texttt{false} & \texttt{\dq{}false\dq{}} & --- \\
\texttt{OM and Size Correct based on target sample} & name & --- & a target sample name, or empty for all \\
\texttt{Source group} & name & --- & a source group name \\
\texttt{Use only selected elements} & \texttt{true} / \texttt{false} & \texttt{\dq{}false\dq{}} & --- \\
\texttt{Use only selected samples} & \texttt{true} / \texttt{false} & \texttt{\dq{}false\dq{}} & --- \\
\end{argtable}

\subsection{\texttt{Source\_Verify}}
\label{cmd:source_verify}
\index{Source\_Verify}

This tool evaluates weather the finger-printing algorithm is capable of correctly infer the source samples based on their elemental profiles. The analysis is done by treating individial source samples from the specified source group as target samples and the use the Levenberg-Marquardt optimization to infer the contribution of sources into it.

\begin{argtable}
\texttt{Apply size and organic matter correction} & \texttt{true} / \texttt{false} & \texttt{\dq{}true\dq{}} & --- \\
\texttt{Softmax transformation} & \texttt{true} / \texttt{false} & \texttt{\dq{}true\dq{}} & --- \\
\texttt{Source Group} & name & --- & a source group name \\
\end{argtable}

\subsection{\texttt{Test CMB Bayesian}}
\label{cmd:test-cmb-bayesian}
\index{Test CMB Bayesian}

\begin{argtable}
\texttt{Number of chains} & whole number & \texttt{\dq{}8\dq{}} & --- \\
\texttt{Number of samples} & whole number & \texttt{\dq{}1000\dq{}} & --- \\
\texttt{Samples to be discarded (burnout)} & whole number & \texttt{\dq{}100\dq{}} & --- \\
\texttt{samples\_file\_name} & file name & \texttt{\dq{}MCMC\_Output.txt\dq{}} & --- \\
\end{argtable}



%======================================================================
\chapter{The report}
\label{ch:report}

A run produces one JSON file describing everything that happened. It is both
the result and the record: it states which project was used, what each step was
asked to do, what it did, and what came out.

\section{What is in it}

\begin{longtable}{@{}p{0.24\textwidth}p{0.68\textwidth}@{}}
\toprule
\textbf{Key} & \textbf{Meaning} \\
\midrule
\endfirsthead
\texttt{script} & The script that was run. \\
\texttt{project} & The project that was loaded. \\
\texttt{working\_folder} & Where analyses wrote files of their own. \\
\texttt{started}, \texttt{finished} & When the run began and ended. \\
\texttt{succeeded} & Whether every step succeeded. \\
\texttt{steps} & One entry per step, in the order they ran. \\
\texttt{stopped\_at\_step} & Present only if the run was abandoned after a
  failure. \\
\bottomrule
\end{longtable}

And within each step:

\begin{longtable}{@{}p{0.24\textwidth}p{0.68\textwidth}@{}}
\toprule
\textbf{Key} & \textbf{Meaning} \\
\midrule
\endfirsthead
\texttt{step}, \texttt{command} & Position in the script, and what was run. \\
\texttt{arguments} & The arguments as given. \\
\texttt{defaulted} & Arguments that were not given and were filled in. \\
\texttt{status} & \texttt{succeeded} or \texttt{failed}. \\
\texttt{seconds} & How long it took. \\
\texttt{warnings} & Anything the analysis reported, such as a reason it
  declined to run. \\
\texttt{results} & What it produced. \\
\texttt{changes} & For a setup command, what it altered. \\
\bottomrule
\end{longtable}

\section{Results}

Results are numbered in the order the analysis produced them, as they are in
the interface: \texttt{1:Contributions}, \texttt{2:Modeled Elemental Profile},
and so on. What each contains depends on the analysis; a contribution is a
figure per source, while a credible interval is a low and high value per
source.

\section{Exit status}

The program exits with status 0 only if every step succeeded, and 1 otherwise.
A script can rely on this:

\begin{lstlisting}
sedsat3-cli run.json && echo "all steps succeeded"
\end{lstlisting}

%======================================================================
\chapter{Options}

\begin{longtable}{@{}p{0.30\textwidth}p{0.62\textwidth}@{}}
\toprule
\textbf{Option} & \textbf{Meaning} \\
\midrule
\endfirsthead
\texttt{--output \textit{file}} & Where to write the report, overriding the
  script. \\
\addlinespace
\texttt{--quiet} & Do not print progress or announce each step. Warnings are
  still printed, and the report is unchanged. \\
\addlinespace
\texttt{--resources \textit{dir}} & Where to find
  \texttt{forms\_structures.json}. \\
\addlinespace
\texttt{--list-commands} & List every command and stop. \\
\addlinespace
\texttt{--describe \textit{cmd}} & Describe one command and stop. \\
\addlinespace
\texttt{--help} & Print usage and stop. \\
\bottomrule
\end{longtable}

%======================================================================
\chapter{Worked examples}

\section{A deterministic workflow}

The sequence the interface presents under \emph{Deterministic Sediment
Fingerprinting}, written down. Note that the steps build on one another: the
correction has to happen before the fingerprinting that relies on it.

\begin{lstlisting}[language=scriptjson]
{
  "project": "my-project.cmb",
  "output": "deterministic.json",
  "on_error": "stop",
  "steps": [
    { "command": "Outlier",
      "arguments": { "Threshold": "3" } },

    { "command": "MLR",
      "arguments": { "Organic Matter constituent": "TOC",
                     "Particle Size constituent": "D50" } },

    { "command": "Bracketing Analysis Batch",
      "arguments": {} },

    { "command": "SDFAM",
      "arguments": { "Use only selected elements": "true" } },

    { "command": "Levenberg-Marquardt-Batch",
      "arguments": { "Apply size and organic matter correction": "true" } }
  ]
}
\end{lstlisting}

\section{A Bayesian run}

\begin{lstlisting}[language=scriptjson]
{
  "project": "my-project.cmb",
  "output": "bayesian.json",
  "steps": [
    { "command": "CMB Bayesian-Batch",
      "arguments": {
        "Number of samples": "1000",
        "Number of chains": "8",
        "Samples to be discarded (burnout)": "100",
        "Apply size and organic matter correction": "false"
      } }
  ]
}
\end{lstlisting}

This writes a good deal besides the report: for each target sample, a folder
containing the chains, the posterior distributions, the credible intervals and
the predicted concentrations. They go under \texttt{working\_folder}.

\section{Comparing two tracer sets}

The case the command line build is best at. Two scripts differing in one line,
each writing its own report, and the two compared afterwards.

\begin{lstlisting}[language=scriptjson]
{
  "project": "my-project.cmb",
  "output": "with-isotopes.json",
  "steps": [
    { "command": "set-elements",
      "arguments": { "only": ["Aluminum","Barium","Lead","Zinc",
                              "TOC","N","C13","N15"] } },
    { "command": "Levenberg-Marquardt-Batch",
      "arguments": { "Apply size and organic matter correction": "false" } }
  ]
}
\end{lstlisting}

and the same with the isotopes left out:

\begin{lstlisting}[language=scriptjson]
    { "command": "set-elements",
      "arguments": { "only": ["Aluminum","Barium","Lead","Zinc"] } },
\end{lstlisting}

Because the only difference between the two runs is recorded in the scripts,
the comparison means something, and it can be repeated.

%======================================================================
\chapter{When something goes wrong}
\label{ch:troubleshooting}

\section{The script was not run}

Every step is checked first, so this means a command or argument name was not
recognised. The message names the step and suggests the closest accepted name.
Names must match the dialogs exactly, including spaces and capitalisation;
\texttt{--describe} is the authority.

\section{A step failed}

The step's entry in the report says so, and its \texttt{warnings} usually say
why. The common ones:

\begin{sloppypar}
\begin{description}[style=nextline,leftmargin=0pt]
  \item[Perform Organic Matter and Size Correction first] An analysis was asked
    to apply the correction, but the correction has not been carried out. Add
    an \texttt{MLR} step before it, or set
    \texttt{\dq{}Apply size and organic matter correction\dq{}} to
    \texttt{\dq{}false\dq{}}.
  \item[Zero or negative values for element \ldots] A constituent has values a
    log-normal distribution cannot represent. Exclude it with
    \texttt{set-elements}, or exclude the samples responsible with
    \texttt{set-samples}.
  \item[Singular matrix in within group scatter matrix] A discriminant analysis
    was given constituents that do not carry independent information. Reduce
    the set.
\end{description}
\end{sloppypar}

\section{Nothing seems to be included in the analysis}

Projects saved by older versions do not record which constituents are included,
and they load with none of them selected. Put a \texttt{set-elements} step at
the start of the script to say explicitly what should be used.

\section{Warnings about a singular system}

Messages of the form

\begin{lstlisting}
warning: solve(): system is singular (rcond: 1.0e-19);
attempting approx solution
\end{lstlisting}

come from the matrix library during the Levenberg-Marquardt fit. They are
normally emitted in the last few iterations, after the fit has converged and
the remaining steps are negligible, and in that case the result is sound ---
check that the contributions are sensible and sum to one.

If they appear from the beginning of a fit rather than at the end, the
parameters are not identifiable from the data: there are too few informative
tracers for the number of sources, or the tracers chosen do not distinguish
them. Reducing the source count or choosing better tracers is the remedy, not
running longer.

\section{It cannot find the command definitions}

The program warns at startup if it cannot find
\texttt{resources/forms\_structures.json}, and the message lists everywhere it
looked. Set \texttt{SEDSAT3\_RESOURCES} or pass \texttt{--resources}. Without
it, scripts are not checked and omitted arguments are not filled in, so
analyses may fail in confusing ways.

\printindex

\end{document}
